\documentclass[11pt,a4paper]{article}
\usepackage[utf8]{inputenc}
\usepackage[T1]{fontenc}
\usepackage[brazilian]{babel}
\usepackage{lmodern}
\usepackage{amsmath,amssymb,amsthm}
\usepackage[margin=2.5cm]{geometry}
\usepackage{algorithm}
\usepackage[noend]{algpseudocode}

\floatname{algorithm}{Algoritmo}
\renewcommand{\algorithmicrequire}{\textbf{Entrada:}}
\renewcommand{\algorithmicensure}{\textbf{Saída:}}

\newtheorem{lema}{Lema}
\newtheorem{teorema}{Teorema}

\renewcommand{\thesubsection}{(\alph{subsection})}

\title{Projeto e Análise de Algoritmos \\ Trabalho 4}
\author{Gideão Santana}
\date{}

\begin{document}
\maketitle

%====================================================================
\section*{Questão 1 -- Empacotamento de livros em envelopes (Algoritmo Guloso)}
\setcounter{lema}{0}\setcounter{teorema}{0}
\noindent\textbf{Notação.} Seja $S$ o conjunto de livros ainda não empacotados,
$C = 5$ a capacidade de um envelope, $h \in S$ um livro de \emph{maior} peso em $S$
e $\ell \in S\setminus\{h\}$ um livro de \emph{menor} peso em $S\setminus\{h\}$.
Um envelope é \emph{viável} se contém no máximo dois livros e peso total $\le C$.
Denotamos por $\mathrm{OPT}(S)$ o número mínimo de envelopes para $S$.

%--------------------------------------------------------------------
\subsection{Estratégia gulosa}

Ordene os livros em ordem não crescente de peso. A cada passo, tome o livro
mais pesado ainda não empacotado, $h$:
\begin{itemize}
  \item se $p_h + p_\ell \le 5$, coloque $h$ e o livro mais leve restante $\ell$
        no mesmo envelope, $\{h,\ell\}$;
  \item caso contrário (ou se $h$ for o único livro restante), coloque $h$ sozinho, $\{h\}$.
\end{itemize}

\noindent\textbf{Por que é promissora.} O livro mais pesado é o mais difícil de
acomodar: se ele não cabe com o livro mais leve, não cabe com nenhum outro, e
portanto ficará sozinho em \emph{qualquer} solução. Se cabe, usar o mais leve como
parceiro é a escolha que menos ``gasta'' livros úteis: os livros de peso
intermediário ficam disponíveis para serem pareados entre si ou com outros
livros pesados.

\noindent\textbf{Exemplo.} Ordenando: $5\,(4.0), 3\,(3.0), 4\,(2.5), 2\,(2.0), 1\,(1.0)$.
Obtemos $\{5,1\}$ ($4+1=5$), $\{3,2\}$ ($3+2=5$) e $\{4\}$: 3 envelopes.

%--------------------------------------------------------------------
\subsection{Propriedade de escolha gulosa}

\begin{lema}[Escolha gulosa]
Seja $S$ não vazio, $h$ um livro de maior peso em $S$ e $\ell$ um livro de menor
peso em $S \setminus \{h\}$. Seja $E$ o envelope escolhido pelo algoritmo, isto é,
$E=\{h,\ell\}$ se $|S|\ge 2$ e $p_h+p_\ell\le 5$, e $E=\{h\}$ caso contrário.
Então existe uma solução ótima para $S$ que contém o envelope $E$.
\end{lema}

\begin{proof}
Seja $O$ uma solução ótima qualquer para $S$.

\smallskip\noindent\emph{Caso (ii): $E = \{h\}$.} Se $|S|=1$, é trivial. Senão,
$p_h + p_\ell > 5$, e para todo $x \in S\setminus\{h\}$ temos $p_x \ge p_\ell$, logo
$p_h + p_x \ge p_h + p_\ell > 5$. Assim nenhum livro pode acompanhar $h$, e em
$O$ (como em qualquer solução) $h$ está sozinho: $\{h\} \in O$.

\smallskip\noindent\emph{Caso (i): $E = \{h,\ell\}$, com $p_h+p_\ell \le 5$.}
Se $\{h,\ell\}\in O$, nada a fazer. Caso contrário, construímos $O'$ por troca:
\begin{enumerate}
  \item \emph{$h$ está sozinho em $O$.} Seja $B$ o envelope de $\ell$ em $O$
    ($B=\{\ell\}$ ou $B=\{\ell,y\}$). Substitua $\{h\}$ e $B$ por $\{h,\ell\}$ e
    $B\setminus\{\ell\}$ (descartando-o se ficar vazio). $\{h,\ell\}$ é viável por
    hipótese e $B\setminus\{\ell\}$ é viável por ser subconjunto de envelope viável.
  \item \emph{$h$ está com $x \ne \ell$ em $O$, ou seja, $\{h,x\} \in O$.}
    Seja $B$ o envelope de $\ell$.
    \begin{itemize}
      \item Se $B=\{\ell\}$: troque $\{h,x\},\{\ell\}$ por $\{h,\ell\},\{x\}$.
      \item Se $B=\{\ell,y\}$: troque $\{h,x\},\{\ell,y\}$ por $\{h,\ell\},\{x,y\}$.
        O envelope $\{x,y\}$ é viável, pois $p_y \le p_h$ ($h$ é o mais pesado) e,
        portanto, $p_x + p_y \le p_x + p_h \le 5$.
    \end{itemize}
\end{enumerate}
Em todos os casos $O'$ é uma partição viável de $S$ com $|O'| \le |O|$, logo
$O'$ é ótima e contém $\{h,\ell\}$.
\end{proof}

%--------------------------------------------------------------------
\subsection{Propriedade de subestrutura ótima}

O subproblema resultante após a escolha gulosa $E$ é $S' = S \setminus E$:
\begin{itemize}
  \item caso (i): $S' = S \setminus \{h,\ell\}$ (os livros restantes, exceto o mais pesado e o mais leve);
  \item caso (ii): $S' = S \setminus \{h\}$ (os livros restantes, exceto o mais pesado).
\end{itemize}
Em ambos os casos $S'$ é uma instância do mesmo problema (mesmas restrições).

\begin{lema}[Subestrutura ótima]
Se $O$ é uma solução ótima para $S$ que contém $E$ (existe pelo Lema 1), então
$O \setminus \{E\}$ é uma solução ótima para $S'$. Consequentemente,
$\mathrm{OPT}(S) = 1 + \mathrm{OPT}(S')$, e $\{E\} \cup O^*$ é ótima para $S$ para
qualquer solução ótima $O^*$ de $S'$.
\end{lema}

\begin{proof}
$O\setminus\{E\}$ particiona $S'$ em envelopes viáveis, com $|O|-1$ envelopes.
Suponha que exista uma solução $O''$ para $S'$ com $|O''| < |O|-1$. Então
$O''\cup\{E\}$ é uma solução viável para $S$ (pois $E$ é viável e disjunto de $S'$)
com $|O''|+1 < |O|$ envelopes, contradizendo a otimalidade de $O$. Logo
$\mathrm{OPT}(S') = |O|-1 = \mathrm{OPT}(S)-1$, e para qualquer $O^*$ ótima de $S'$,
$|\{E\}\cup O^*| = 1+\mathrm{OPT}(S') = \mathrm{OPT}(S)$.
\end{proof}

%--------------------------------------------------------------------
\subsection{Pseudocódigo}

Após ordenar em ordem não crescente, o índice $i$ aponta para o livro mais pesado
restante e $j$ para o mais leve restante (técnica de dois ponteiros).

\begin{algorithm}[H]
\caption{\textsc{Empacota-Livros}$(p, n)$}
\begin{algorithmic}[1]
\Require pesos $p_1,\dots,p_n$, com $0 < p_k \le 5$
\Ensure lista $\mathcal{E}$ de envelopes com o menor número possível de envelopes
\State $L \gets$ livros $1..n$ ordenados em ordem \textbf{não crescente} de peso \Comment{$O(n\lg n)$}
\State $\mathcal{E} \gets \emptyset$
\State $i \gets 1$; \quad $j \gets n$
\While{$i \le j$}
  \If{$i < j$ \textbf{and} $p_{L[i]} + p_{L[j]} \le 5$} \Comment{caso (i)}
    \State $\mathcal{E} \gets \mathcal{E} \cup \{\{L[i], L[j]\}\}$
    \State $j \gets j - 1$
  \Else \Comment{caso (ii) ou último livro}
    \State $\mathcal{E} \gets \mathcal{E} \cup \{\{L[i]\}\}$
  \EndIf
  \State $i \gets i + 1$
\EndWhile
\State \Return $\mathcal{E}$
\end{algorithmic}
\end{algorithm}

Invariante: no início de cada iteração, os livros não empacotados são exatamente
$L[i..j]$, com $L[i]$ o mais pesado e $L[j]$ o mais leve entre eles.

%--------------------------------------------------------------------
\subsection{Correção}

\begin{teorema}
\textsc{Empacota-Livros} devolve uma distribuição viável com $\mathrm{OPT}(S)$ envelopes.
\end{teorema}

\begin{proof}
Por indução em $|S|$ (o número de livros em $L[i..j]$).
\emph{Base:} $|S| = 0$: o algoritmo devolve $\emptyset$, que é ótimo.
\emph{Passo:} seja $|S|\ge 1$. Pelo invariante, o algoritmo escolhe exatamente o
envelope $E$ dos itens (a)--(b): $\{L[i],L[j]\}$ se cabem juntos e $i<j$, ou
$\{L[i]\}$ caso contrário; $E$ é viável. Pelo Lema~1, existe solução ótima de $S$
que contém $E$. O algoritmo prossegue sobre $S' = S\setminus E = L[i'..j']$
(subproblema identificado em (c)), com $|S'| < |S|$; pela hipótese de indução,
devolve uma solução ótima $O^*$ de $S'$. Pelo Lema~2, $\{E\}\cup O^*$ é ótima para
$S$, com $1+\mathrm{OPT}(S') = \mathrm{OPT}(S)$ envelopes. Como cada livro entra em
exatamente um envelope (cada índice sai de $L[i..j]$ uma única vez), a saída é uma
partição válida.
\end{proof}

%--------------------------------------------------------------------
\subsection{Complexidade}

\begin{itemize}
  \item \textbf{Ordenação:} $O(n \lg n)$ com um algoritmo ótimo por comparação
        (\textit{Merge Sort} ou \textit{Heapsort}).
  \item \textbf{Laço:} a cada iteração $i$ aumenta em 1 (e $j$ possivelmente
        diminui em 1), logo $j-i$ decresce ao menos 1 por iteração e há no máximo
        $n$ iterações, cada uma em $O(1)$: total $O(n)$.
  \item \textbf{Construção da saída:} cada envelope é inserido em $O(1)$; são no
        máximo $n$ envelopes: $O(n)$.
\end{itemize}
Portanto, $T(n) = O(n\lg n) + O(n) + O(n) = O(n \lg n)$, dominado pela ordenação.
O espaço adicional é $O(n)$.


\section*{Questão 2 -- Paradas em hotéis (Algoritmo Guloso)}
\setcounter{lema}{0}\setcounter{teorema}{0}
\noindent\textbf{Notação.} Uma sequência de hotéis $S=[s_1,\dots,s_r]$ é
\emph{viável} se $s_1 = 1$, $s_r = n$, $s_i < s_{i+1}$ e
$x[s_{i+1}] - x[s_i] \le D$ para todo $1 \le i < r$. A escolha gulosa a partir
da parada $u$ é $\max\{\, w > u : x[w] - x[u] \le D \,\}$, que existe para
$u<n$, pois $x[u+1]-x[u]\le D$.

\noindent\textbf{Exemplo} ($D=12$): a partir do hotel 1 ($x=0$), o mais distante
alcançável é o hotel 4 ($x=11$); a partir do hotel 4, é o hotel 6 ($x=17$).
Logo $G=[1,4,6]$.

%--------------------------------------------------------------------
\subsection{Propriedade de escolha gulosa (argumento de troca)}

Sejam $G=[u_1,\dots,u_m]$ e $O=[v_1,\dots,v_\ell]$ como no enunciado, com
$v_i = u_i$ para $1\le i<k$, $v_k \ne u_k$ e $1<k<\ell$.

\begin{lema}
$v_k < u_k$.
\end{lema}
\begin{proof}
Como $O$ é viável e $v_{k-1}=u_{k-1}$, temos $x[v_k] - x[u_{k-1}] \le D$ e
$v_k > u_{k-1}$. Logo $v_k$ é um candidato da escolha gulosa a partir de $u_{k-1}$,
e $u_k$ é o maior candidato: $v_k \le u_k$. Como $v_k \ne u_k$, $v_k < u_k$.
\end{proof}

\noindent\textbf{Construção.} Troque $v_k$ por $u_k$ em $O$, descartando os
hotéis de $O$ que ficaram ``para trás'' de $u_k$:
\[
O' = [\,v_1,\dots,v_{k-1},\; u_k,\; v_{j+1},\dots,v_\ell\,],
\qquad j = \max\{\, i \ge k : v_i \le u_k \,\}.
\]
(Note que $j\ge k$, pois $v_k<u_k$. Se $v_j = u_k$, então $v_j$ e $u_k$ são o mesmo
hotel e ele aparece uma única vez.)

\begin{lema}
$O'$ é viável e $|O'| \le |O|$.
\end{lema}
\begin{proof}
\emph{Tamanho:} $O'$ remove de $O$ os $j-k+1 \ge 1$ hotéis $v_k,\dots,v_j$ e
insere apenas $u_k$, logo $|O'| = \ell - (j-k) \le \ell$.

\emph{Ordem:} $v_{k-1}=u_{k-1}<u_k$ e, pela escolha de $j$, $u_k < v_{j+1}$.

\emph{Distâncias:} os trechos $v_1\to\cdots\to v_{k-1}$ e
$v_{j+1}\to\cdots\to v_\ell$ são trechos de $O$, portanto $\le D$. Restam dois trechos:
\begin{itemize}
  \item $u_{k-1}\to u_k$: $x[u_k]-x[u_{k-1}] \le D$ pela definição da escolha gulosa;
  \item $u_k \to v_{j+1}$ (se $j<\ell$): como $x[v_j] \le x[u_k]$,
        \[ x[v_{j+1}] - x[u_k] \;\le\; x[v_{j+1}] - x[v_j] \;\le\; D. \]
\end{itemize}
\emph{Extremos:} $O'$ começa em $v_1=1$. Se $j<\ell$, termina em $v_\ell=n$;
se $j=\ell$, então $n = v_\ell \le u_k \le n$, logo $u_k=n$ e $O'$ termina em $n$.
\end{proof}

Como $O$ é ótima e $|O'|\le|O|$, $O'$ também é ótima (logo, na verdade, $j=k$ e
$|O'|=|O|$), e $O'$ coincide com $G$ nas $k$ primeiras paradas. Isso prova a
sentença: é possível trocar $v_k$ por $u_k$ sem piorar a solução ótima.

\begin{teorema}
A solução gulosa $G$ é ótima, ou seja, $m = \ell$.
\end{teorema}
\begin{proof}
Aplicando a troca repetidamente, o primeiro índice de divergência cresce a cada
passo; obtemos assim uma solução ótima $O^*=[v_1,\dots,v_\ell]$ com $v_i=u_i$
para todo $1\le i\le \ell-1$. Como $O^*$ é viável, $x[n]-x[u_{\ell-1}] =
x[v_\ell]-x[v_{\ell-1}] \le D$; então o hotel $n$ (o mais distante de todos) é
alcançável a partir de $u_{\ell-1}$ e a escolha gulosa é $u_\ell = n$. Assim o
guloso termina com $m=\ell$ paradas. (Equivalentemente, pelo subproblema descrito no
enunciado, a cada troca o problema restante parte de $x[u_k]$ e o argumento se
repete por indução.)
\end{proof}

%--------------------------------------------------------------------
\subsection{Pseudocódigo e complexidade}

\begin{algorithm}[H]
\caption{\textsc{Paradas-Hotéis}$(x, n, D)$}
\begin{algorithmic}[1]
\Require posições $x[1..n]$ em ordem crescente, com $x[i+1]-x[i]\le D$
\Ensure lista $G$ com o menor número de paradas, de $1$ até $n$
\State $G \gets [1]$
\State $u \gets 1$ \Comment{última parada}
\State $i \gets 1$ \Comment{candidato atual}
\While{$u < n$}
  \While{$i < n$ \textbf{and} $x[i+1] - x[u] \le D$}
    \State $i \gets i + 1$ \Comment{avança enquanto o próximo hotel é alcançável}
  \EndWhile
  \State $u \gets i$ \Comment{escolha gulosa: hotel mais distante alcançável}
  \State \textsc{Anexa}$(G, u)$
\EndWhile
\State \Return $G$
\end{algorithmic}
\end{algorithm}

\noindent\textbf{Terminação e corretude do laço.} No início de cada iteração do
laço externo, $i = u$. O laço interno avança $i$ ao maior índice com
$x[i]-x[u]\le D$ (a escolha gulosa); como $x[u+1]-x[u]\le D$, ele avança pelo menos
uma vez, logo $u$ cresce estritamente e o algoritmo termina com $u=n$.

\noindent\textbf{Complexidade.} O índice $i$ nunca decresce e vai de $1$ até $n$,
então o laço interno executa no total no máximo $n-1$ incrementos ao longo de toda
a execução. O laço externo executa no máximo $n-1$ vezes, cada uma com custo
$O(1)$ fora do laço interno (anexar a $G$ custa $O(1)$). Logo o tempo total é
\[ T(n) = O(n), \]
sem necessidade de ordenação, pois $x$ já está em ordem crescente. O espaço adicional
é $O(n)$ para armazenar $G$ (ou $O(1)$, se as paradas forem apenas impressas).

\clearpage

%====================================================================
\section*{Questão 3 -- Impressão nítida (Programação Dinâmica)}
\setcounter{lema}{0}\setcounter{teorema}{0}
\noindent\textbf{Ideia.} No algoritmo recursivo, a única informação que muda entre
as chamadas é $i$, o índice da primeira palavra do subproblema ($M$ e $W$ são
fixos). Como o mesmo $i$ é resolvido muitas vezes por caminhos diferentes de quebras,
memorizamos a resposta de cada $i$ em vetores:
\begin{itemize}
  \item $\mathit{custo}[i]$: custo mínimo para imprimir as palavras $i,\dots,n$;
  \item $\mathit{prox}[i]$: índice $j$ da próxima quebra de linha em uma solução
        ótima a partir de $i$ ($\mathit{prox}[i]=n+1$ indica que não há quebra).
\end{itemize}
A recorrência, tirada diretamente das linhas 01--11 do algoritmo recursivo, é
\[
\mathit{custo}[i] =
\begin{cases}
0, & \text{se } M[i,n] \le W,\\[2pt]
\displaystyle\min_{\substack{i < j \le n\\ M[i,j-1]\le W}}
  \Big\{ (W - M[i,j-1])^2 + \mathit{custo}[j] \Big\}, & \text{caso contrário.}
\end{cases}
\]
O conjunto do $\min$ nunca é vazio, pois $M[i,i]=L[i]\le W$ (o candidato $j=i+1$
sempre é válido). A resposta do problema é $\mathit{custo}[1]$.

%--------------------------------------------------------------------
\subsection{Número de subsoluções}

São armazenadas $n$ subsoluções: uma para cada $i \in \{1,\dots,n\}$, isto é,
$\mathit{custo}[1..n]$ (e o respectivo $\mathit{prox}[1..n]$).
\emph{Justificativa:} todo subproblema gerado pela recursão é ``imprimir as palavras
$i,\dots,n$'', identificado unicamente pelo parâmetro $i$, que só assume valores entre
$1$ e $n$ (a chamada inicial usa $i=1$ e as recursivas usam $j\in\{i+1,\dots,n\}$).
Portanto há apenas $n$ subproblemas distintos, embora a recursão sem memorização os
resolva um número exponencial de vezes.

%--------------------------------------------------------------------
\subsection{Caso base}

O caso base corresponde às linhas 01--02: para todo $i$ tal que $M[i,n] \le W$, as
palavras $i,\dots,n$ cabem em uma única linha, que é a última e tem custo zero.
Assim,
\[
\mathit{custo}[i] = 0 \quad\text{e}\quad \mathit{prox}[i] = n+1
\qquad \text{para todo } i \text{ com } M[i,n]\le W.
\]
Como $M[i,n]$ decresce quando $i$ cresce, esses índices formam um sufixo
$\{i_0,\dots,n\}$; em particular, $i=n$ é sempre caso base, pois $M[n,n]=L[n]\le W$.

%--------------------------------------------------------------------
\subsection{Ordem de preenchimento}

A tabela é preenchida com $i$ \textbf{decrescente}, de $n$ até $1$.

\emph{Argumento:} $\mathit{custo}[i]$ depende apenas de $\mathit{custo}[j]$ com
$j > i$ (na recorrência, $j$ varia de $i+1$ a $n$). Quando calculamos a posição $i$,
todas as posições $i+1,\dots,n$ já foram preenchidas, seja pelo caso base, seja
em iterações anteriores. Logo nenhum elemento ainda não preenchido é acessado.
A matriz $M$ é preenchida antes, pelo algoritmo do enunciado.

%--------------------------------------------------------------------
\subsection{Pseudocódigo}

\begin{algorithm}[H]
\caption{\textsc{Impressão-Nítida-PD}$(L, n, W)$}
\begin{algorithmic}[1]
\Require tamanhos $L[1..n]$, com $L[i]\le W$, e largura $W$
\Ensure custo mínimo e conjunto de quebras de uma solução ótima
\For{$i \gets 1$ \textbf{até} $n$} \Comment{preenche $M$ em $O(n^2)$}
  \State $M[i,i] \gets L[i]$
  \For{$j \gets i+1$ \textbf{até} $n$}
    \State $M[i,j] \gets M[i,j-1] + L[j] + 1$
  \EndFor
\EndFor
\For{$i \gets n$ \textbf{até} $1$ \textbf{passo} $-1$} \Comment{ordem decrescente}
  \If{$M[i,n] \le W$} \Comment{caso base: última linha}
    \State $\mathit{custo}[i] \gets 0$; \quad $\mathit{prox}[i] \gets n+1$
  \Else
    \State $\mathit{custo}[i] \gets \infty$
    \For{$j \gets i+1$ \textbf{até} $n$}
      \If{$M[i,j-1] > W$} \textbf{break} \EndIf
      \State $c \gets (W - M[i,j-1])^2 + \mathit{custo}[j]$
      \If{$c < \mathit{custo}[i]$}
        \State $\mathit{custo}[i] \gets c$; \quad $\mathit{prox}[i] \gets j$
      \EndIf
    \EndFor
  \EndIf
\EndFor
\State $\mathit{sol} \gets \emptyset$; \quad $i \gets 1$ \Comment{reconstrução da solução}
\While{$\mathit{prox}[i] \le n$}
  \State $i \gets \mathit{prox}[i]$
  \State $\mathit{sol} \gets \mathit{sol} \cup \{i\}$
\EndWhile
\State \Return $(\mathit{sol}, \mathit{custo}[1])$
\end{algorithmic}
\end{algorithm}

\noindent\textbf{Exemplo.} Para a entrada do enunciado ($W=12$), o algoritmo
obtém $\mathit{custo}[1]=33$ com quebras em $\{3,5,8,11\}$. Logo, a solução
apresentada no enunciado é ótima.

%--------------------------------------------------------------------
\subsection{Complexidade de tempo}

\begin{itemize}
  \item \textbf{Preenchimento de $M$:} $\sum_{i=1}^{n}(n-i+1) = O(n^2)$.
  \item \textbf{Preenchimento de $\mathit{custo}$ e $\mathit{prox}$:} são $n$
        subproblemas (item (a)); cada um examina no máximo $n-i$ candidatos $j$,
        cada um em $O(1)$ (a linha 10 do recursivo vira a atribuição
        $\mathit{prox}[i]\gets j$, de custo constante). Total:
        $\sum_{i=1}^{n} O(n-i) = O(n^2)$.
  \item \textbf{Reconstrução:} cada passo avança $i$ estritamente, logo $O(n)$.
\end{itemize}
Portanto, o tempo total é
\[ T(n) = O(n^2) + O(n^2) + O(n) = O(n^2), \]
polinomial, em contraste com o tempo exponencial do backtracking. O espaço é
$O(n^2)$ por causa de $M$ (os vetores $\mathit{custo}$ e $\mathit{prox}$ ocupam $O(n)$).

\end{document}
